\chapter{Hematocrit}

\section*{What Is This Test?}
The hematocrit (HCT), also known as packed cell volume (PCV), measures the proportion of whole blood occupied by red blood cells. It is expressed as a percentage and is a key component of the complete blood count (CBC). The hematocrit reflects both the number and size of circulating red blood cells. Because red blood cells are responsible for oxygen delivery to tissues, the hematocrit provides critical information about the blood's oxygen-carrying capacity. A low hematocrit is the primary quantitative indicator of anemia \cite{who2011haemoglobin}, while an elevated hematocrit suggests polycythemia or hemoconcentration. The hematocrit is used in conjunction with hemoglobin concentration and red blood cell count to assess red blood cell mass and classify disorders of red blood cell production, destruction, or loss. It is also used to estimate blood loss during hemorrhage and to monitor response to blood transfusion, erythropoietin therapy, and chemotherapy.

\section*{Alternative Names}
HCT; PCV (Packed Cell Volume); Crit; Hematocrit (Blood); Hct Level

\section*{Principle}
In modern automated hematology analyzers, the hematocrit is typically calculated rather than directly measured \cite{tietz2023textbook}. The analyzer determines the mean corpuscular volume (MCV) by impedance measurement and the red blood cell count (RBC), then calculates hematocrit using the formula: HCT = MCV $\times$ RBC / 10. In older methodology, the manual spun hematocrit involves centrifuging a microcapillary tube of whole blood at 10,000--15,000 rpm for 5 minutes, after which the height of the packed red cell column is measured against the total column height and expressed as a percentage \cite{bain2017dacie}. The spun hematocrit yields a slightly higher value than the automated calculated method (by approximately 1--3\%) because it includes trapped plasma between packed cells. Automated impedance-based analyzers (Sysmex XN, Beckman Coulter DxH, Siemens ADVIA) provide precise and reproducible hematocrit values as part of the routine CBC.

\section*{History}
The concept of measuring packed red cell volume was first described by Magnus Blix in 1890, and the term ``hematocrit'' was coined by Swedish physiologist Hedin in 1891. The microhematocrit centrifuge method was developed in the 1940s and became the standard method for decades, using a high-speed centrifuge and glass capillary tubes. Wintrobe's classic text ``Clinical Hematology'' (1942) established the packed cell volume as a fundamental test in anemia evaluation \cite{kaushansky2021williams}. The introduction of the Coulter Counter in the 1950s enabled automated calculation of hematocrit from directly measured MCV and RBC count parameters. By the 1980s, automated hematology analyzers had largely replaced manual spun hematocrit in clinical laboratories, though the spun method remains useful in resource-limited settings and for point-of-care testing.

\section*{Why This Test Is Ordered}
\begin{itemize}
  \item Screening and diagnosis of anemia --- when hemoglobin and hematocrit are below normal, anemia is present and must be classified by etiology
  \item Evaluation of polycythemia --- elevated hematocrit above 48\% (women) or 52\% (men) suggests polycythemia vera, secondary polycythemia from chronic hypoxia, or relative polycythemia from dehydration
  \item Monitoring of blood loss --- serial hematocrit measurements track the severity of acute or chronic hemorrhage and guide transfusion decisions
  \item Assessment of hydration status --- elevated hematocrit may indicate dehydration and hemoconcentration
  \item Monitoring response to erythropoietin therapy in patients with chronic kidney disease on dialysis
  \item Monitoring hematologic toxicity of chemotherapy and bone marrow recovery following stem cell transplantation
  \item Preoperative assessment and post-surgical monitoring
  \item Evaluation of conditions causing hemolysis, including autoimmune hemolytic anemia, hereditary spherocytosis, and G6PD deficiency
\end{itemize}

\section*{Primary vs Secondary Test}
The hematocrit is a primary component of the complete blood count and serves as a first-line screening test for anemia and polycythemia \cite{tierney2023current}.

\section*{How This Test Is Performed}
A venous blood sample is collected into a lavender-top (EDTA) tube from a peripheral vein, typically the antecubital fossa. The tube is gently inverted 8--10 times to mix the anticoagulant. The sample is processed on an automated hematology analyzer within 4 hours of collection. For a manual spun hematocrit, whole blood is drawn into a heparinized capillary tube (red-tip) from a fingerstick or venous sample, one end is sealed with clay, and the tube is centrifuged in a microhematocrit centrifuge for 5 minutes. The packed cell volume is read against a calibrated scale. The entire venipuncture procedure takes less than 5 minutes.

\section*{What Preparation Is Needed}
No special preparation is required for a routine hematocrit. Fasting is not necessary. Patients should avoid vigorous exercise for 30 minutes prior to sample collection, as this can transiently increase hematocrit through plasma volume shifts and splenic contraction. Excessive tourniquet application during venipuncture (>1 minute) can cause hemoconcentration and falsely elevate results \cite{wallach2022interpretation}. Patients should inform their healthcare provider of current medications, particularly erythropoietin, testosterone, diuretics, and recent blood transfusion.

\section*{Contraindications and Precautions}
\begin{itemize}
  \item No absolute contraindications. Factors affecting results include: prolonged tourniquet application (stasis causes hemoconcentration and falsely elevated hematocrit), recent vigorous exercise, dehydration, pregnancy (physiological hemodilution causes decreased hematocrit, reaching a nadir at 30--34 weeks gestation), living at high altitude (>1,500 meters causes compensatory erythrocytosis), smoking (carbon monoxide-induced polycythemia), immediate post-transfusion measurement, and severe lipemia (may falsely elevate hematocrit on some automated analyzers). EDTA (purple-top tube) is the standard anticoagulant. Samples should be analyzed within 4 hours at room temperature or 24 hours if refrigerated at 2--8\textdegree{}C.
\end{itemize}
\section*{Risks}
Minimal risks are associated with venipuncture: slight pain or stinging at the needle insertion site, bruising (ecchymosis) at the venipuncture site, hematoma formation if inadequate pressure is applied after needle withdrawal, and rarely, vasovagal syncope (fainting). Infection at the puncture site is exceedingly rare.

\section*{What Happens After the Test}
Results are typically available within 30--60 minutes when processed by automated analyzer. The healthcare provider reviews the hematocrit in the context of the full CBC and clinical presentation. If anemia is identified, the MCV, RDW, reticulocyte count, and peripheral blood smear are reviewed to classify the anemia type. If polycythemia is identified, additional testing for JAK2 V617F mutation and erythropoietin level may be ordered to differentiate primary from secondary causes. Serial monitoring may be ordered for hospitalized patients with ongoing blood loss or hemodilution.

\section*{Understanding Results}

\subsection*{Below Normal}
\begin{itemize}
  \item \textbf{Iron deficiency anemia:} Microcytic hypochromic anemia with decreased hematocrit, low MCV, elevated RDW, low ferritin, and low transferrin saturation. The most common cause of low hematocrit worldwide.
  \item \textbf{Anemia of chronic disease:} Normocytic normochromic anemia (can be microcytic in advanced stages) associated with chronic infections (tuberculosis, osteomyelitis), chronic inflammatory conditions (rheumatoid arthritis, systemic lupus erythematosus), or malignancy.
  \item \textbf{Thalassemia:} Microcytic anemia with normal or elevated RDW, caused by defective globin chain synthesis. Alpha-thalassemia and beta-thalassemia trait produce mild anemia; thalassemia major produces severe anemia requiring transfusion.
  \item \textbf{Vitamin B12 or folate deficiency:} Macrocytic anemia (elevated MCV >100 fL) with low hematocrit caused by impaired DNA synthesis, seen in pernicious anemia, gastric bypass, chronic alcoholism, and dietary deficiency.
  \item \textbf{Hemolytic anemia:} Decreased hematocrit due to premature red blood cell destruction. Causes include autoimmune hemolytic anemia (positive direct Coombs test), hereditary spherocytosis, sickle cell disease, G6PD deficiency, and microangiopathic hemolytic anemia (TTP, HUS, DIC).
  \item \textbf{Acute blood loss:} Sudden decrease in hematocrit following trauma, gastrointestinal bleeding, ruptured aortic aneurysm, or postpartum hemorrhage. Initial hematocrit may be normal before hemodilution occurs.
  \item \textbf{Chronic kidney disease:} Normocytic anemia due to decreased erythropoietin production by diseased kidneys, typically when GFR falls below 30 mL/min/1.73m\textsuperscript{2}.
  \item \textbf{Bone marrow failure:} Aplastic anemia, myelodysplastic syndrome, leukemia infiltration, or chemotherapy/radiation-induced myelosuppression producing pancytopenia with low hematocrit.
  \item \textbf{Hemodilution:} Pregnancy (physiological), excessive IV fluid administration, congestive heart failure with fluid overload, or syndrome of inappropriate antidiuretic hormone (SIADH).
\end{itemize}

\subsection*{Above Normal}
\begin{itemize}
  \item \textbf{Polycythemia vera:} Primary myeloproliferative neoplasm with JAK2 V617F mutation in >95\% of cases, characterized by unregulated red blood cell production and elevated hematocrit (often >55\% in men, >50\% in women), often accompanied by elevated WBC and platelet counts.
  \item \textbf{Secondary polycythemia (appropriate):} Compensatory increase in erythropoietin due to chronic hypoxia --- chronic obstructive pulmonary disease (COPD), sleep apnea, congenital cyanotic heart disease, high-altitude living, chronic carbon monoxide exposure.
  \item \textbf{Secondary polycythemia (inappropriate):} Pathologic erythropoietin overproduction from renal cell carcinoma, hepatocellular carcinoma, cerebellar hemangioblastoma, pheochromocytoma, or uterine leiomyoma.
  \item \textbf{Relative polycythemia (hemoconcentration):} Dehydration from vomiting, diarrhea, diuretic use, burns, or heatstroke; Gaisb\"{o}ck syndrome (stress polycythemia with hypertension, obesity, and smoking).
  \item \textbf{Testosterone therapy:} Exogenous testosterone stimulates erythropoiesis, commonly elevating hematocrit in men receiving testosterone replacement therapy.
\end{itemize}

\section*{Normal Results}
\begin{itemize}
  \item \textbf{Adult males:} 38.3--48.6\% \cite{fischbach2023manual}
  \item \textbf{Adult females:} 35.5--44.9\%
  \item \textbf{Children (6--12 years):} 35--45\%
  \item \textbf{Children (1--6 years):} 30--40\%
  \item \textbf{Newborns (first week):} 42--65\% (decreases to 35--49\% by 1 month)
  \item \textbf{Pregnancy (third trimester):} 28--40\% (physiological hemodilution)
  \item \textbf{Adults living at high altitude (>2,400 m):} Males 45--61\%; Females 41--56\%
\end{itemize}

\section*{Follow-up Tests}
\begin{itemize}
  \item \textbf{Complete blood count with differential:} Evaluate hemoglobin, MCV, MCH, MCHC, RDW, WBC, and platelet count to identify patterns of anemia or polycythemia
  \item \textbf{Peripheral blood smear:} Microscopic examination to assess red blood cell morphology --- microcytosis, macrocytosis, hypochromia, target cells, spherocytes, schistocytes, sickle cells, or teardrop cells
  \item \textbf{Reticulocyte count:} Distinguishes between decreased production (low reticulocyte count in iron deficiency, aplastic anemia) and increased destruction/loss (elevated reticulocyte count in hemolysis or hemorrhage)
  \item \textbf{Iron studies:} Serum iron, ferritin, total iron-binding capacity (TIBC), and transferrin saturation to evaluate iron deficiency anemia
  \item \textbf{Vitamin B12 and serum folate levels:} Diagnostic for megaloblastic anemia when MCV is elevated
  \item \textbf{Hemoglobin electrophoresis:} Identifies hemoglobinopathies (sickle cell disease, thalassemia) and quantifies HbA, HbA2, HbF, and HbS
  \item \textbf{Direct antiglobulin test (Coombs test):} Detects autoimmune hemolytic anemia when hemolysis is suspected
  \item \textbf{JAK2 V617F mutation analysis:} Diagnostic for polycythemia vera when hematocrit is persistently elevated
  \item \textbf{Serum erythropoietin level:} Distinguishes primary (low EPO) from secondary (high EPO) polycythemia
  \item \textbf{Bone marrow aspiration and biopsy:} Evaluates bone marrow cellularity and morphology when aplastic anemia, myelodysplastic syndrome, or leukemia is suspected
\end{itemize}

\section*{References}
\bibliographystyle{unsrtnat}
\bibliography{references}

\clearpage
\section*{Regulatory Status and Authorization Date}
This chapter describes a test category, not a specific commercial product. FDA, CDSCO, EU IVDR, and MHRA approval or registration dates therefore cannot be assigned without the assay manufacturer, product name, instrument, and intended use. For laboratory-developed tests, an individual agency approval date may not exist. See the Preface for the interpretation of regulatory status.

\clearpage
